\documentclass[10pt,letterpaper]{article}
\usepackage{cornell-notes}
\usepackage{mathtools}

\title{Chapter 11: Eigensystems}
\author{}
\date{}

\setDocTitle{Chapter 11: Eigensystems}
\setDocSubtitle{Numerical Methods}
\setDocVersion{v1.0}
\setDocStatus{Study Companion}
\setDocOwner{Numerical Methods Cornell Notes}
\setDocAuthor{}
\setDocDate{\today}
\setDocSummary{Cornell-format study and review notes for Chapter 11: Eigensystems.}

\setCornellCollection{Numerical Methods Cornell Notes}
\setCornellUnitType{Chapter}
\setCornellUnitNumber{11}
\setCornellUnitTitle{Eigensystems}
\setCornellCompanionLabel{Study and review companion}

\hypersetup{pdftitle={Chapter 11: Eigensystems},pdfsubject={Cornell notes},pdfauthor={}}

\begin{document}
\maketitle

\section*{Chapter Overview}
This chapter organizes the mathematical and computational ideas behind eigensystems. The notes preserve the numbered section sequence while emphasizing method selection, explicit assumptions, numerical reliability, cost, and verification.

\section*{Learning Objectives}
Explain the governing problem class; compare Jacobi Transformations of a Symmetric Matrix, Real Symmetric Matrices, Reduction of a Symmetric Matrix to Tridiagonal Form: Givens and Householder Reductions, and Eigenvalues and Eigenvectors of a Tridiagonal Matrix; design defensible stopping and verification checks; distinguish problem conditioning from algorithmic stability.

\section*{Prerequisites}
Matrix factorizations, orthogonality, complex numbers, and norms.

\section*{Operational Workflow}
Exploit symmetry or sparsity; balance or scale when appropriate; reduce to structured form; iterate and deflate; compute eigenvectors only as needed; verify normalized residuals and orthogonality.

\subsection*{Formula and notation anchor}
\[
Av=\lambda v,\qquad \rho=\frac{\|Av-\lambda v\|_2}{\|A\|_2\|v\|_2}
\]
The residual measures how accurately a pair satisfies the equation; sensitivity also depends on spectral separation and normality.

\begin{warnbox}{Numerical warning}
Nonnormal sensitivity, clustered eigenvalues, loss of orthogonality, unsafe deflation, incorrect complex-pair handling, and interpreting small residuals without eigenvector conditioning.
\end{warnbox}

\section{11.0 Introduction}
\begin{cornell}
\noterow{What is the central idea?}{Eigenproblems reveal invariant directions and modal scales; symmetry, sparsity, desired spectrum, and conditioning drive algorithm choice.}
\noterow{How should it be used?}{For Introduction, begin from explicit inputs, outputs, preconditions, and representation choices.  Exploit symmetry or sparsity; balance or scale when appropriate; reduce to structured form; iterate and deflate; compute eigenvectors only as needed; verify normalized residuals and orthogonality.}
\noterow{Cost and reliability}{Analyze arithmetic work, storage, data movement, and convergence separately.  Relevant failure pressures include nonnormal sensitivity, clustered eigenvalues, loss of orthogonality, unsafe deflation, incorrect complex-pair handling, and interpreting small residuals without eigenvector conditioning.  Use scaled tolerances and report both success criteria and diagnostic evidence.}
\noterow{Implementation checklist}{\begin{itemize}
\item Validate dimensions, domains, ordering, and structural assumptions.
\item Guard small divisors, invalid states, overflow, underflow, and nonfinite values.
\item Make mutation, workspace, indexing, and termination behavior explicit.
\item Test a simple exact case, a difficult valid case, a boundary case, and an expected failure.
\end{itemize}}
\end{cornell}
\begin{summarybox}{Section summary}
Eigenproblems reveal invariant directions and modal scales; symmetry, sparsity, desired spectrum, and conditioning drive algorithm choice.  Select this technique only when its assumptions match the data, and accept a result only after an independent residual, error estimate, invariant, or refinement check.
\end{summarybox}

\section{11.1 Jacobi Transformations of a Symmetric Matrix}
\begin{cornell}
\noterow{What is the central idea?}{Jacobi rotations annihilate selected off-diagonal entries while preserving symmetry and eigenvalues, accumulating orthogonal eigenvectors.}
\noterow{How should it be used?}{For Jacobi Transformations of a Symmetric Matrix, begin from explicit inputs, outputs, preconditions, and representation choices.  Exploit symmetry or sparsity; balance or scale when appropriate; reduce to structured form; iterate and deflate; compute eigenvectors only as needed; verify normalized residuals and orthogonality.}
\noterow{Quantitative anchor}{\[
A_{k+1}=J_k^T A_k J_k
\]
State the dimensions, scaling, and validity assumptions before using this relation.  Check the resulting residual or invariant rather than trusting an iteration count alone.}
\noterow{Cost and reliability}{Analyze arithmetic work, storage, data movement, and convergence separately.  Relevant failure pressures include nonnormal sensitivity, clustered eigenvalues, loss of orthogonality, unsafe deflation, incorrect complex-pair handling, and interpreting small residuals without eigenvector conditioning.  Use scaled tolerances and report both success criteria and diagnostic evidence.}
\noterow{Implementation checklist}{\begin{itemize}
\item Validate dimensions, domains, ordering, and structural assumptions.
\item Guard small divisors, invalid states, overflow, underflow, and nonfinite values.
\item Make mutation, workspace, indexing, and termination behavior explicit.
\item Test a simple exact case, a difficult valid case, a boundary case, and an expected failure.
\end{itemize}}
\end{cornell}
\begin{summarybox}{Section summary}
Jacobi rotations annihilate selected off-diagonal entries while preserving symmetry and eigenvalues, accumulating orthogonal eigenvectors.  Select this technique only when its assumptions match the data, and accept a result only after an independent residual, error estimate, invariant, or refinement check.
\end{summarybox}

\section{11.2 Real Symmetric Matrices}
\begin{cornell}
\noterow{What is the central idea?}{Real symmetry guarantees real eigenvalues and an orthonormal eigenbasis, enabling stable reductions and strong residual checks.}
\noterow{How should it be used?}{For Real Symmetric Matrices, begin from explicit inputs, outputs, preconditions, and representation choices.  Exploit symmetry or sparsity; balance or scale when appropriate; reduce to structured form; iterate and deflate; compute eigenvectors only as needed; verify normalized residuals and orthogonality.}
\noterow{Quantitative lens}{Identify the governing equation, recurrence, probability law, or geometric predicate for Real Symmetric Matrices.  Define every variable and scale; then derive a residual, error estimate, or invariant that can be checked independently.}
\noterow{Cost and reliability}{Analyze arithmetic work, storage, data movement, and convergence separately.  Relevant failure pressures include nonnormal sensitivity, clustered eigenvalues, loss of orthogonality, unsafe deflation, incorrect complex-pair handling, and interpreting small residuals without eigenvector conditioning.  Use scaled tolerances and report both success criteria and diagnostic evidence.}
\noterow{Implementation checklist}{\begin{itemize}
\item Validate dimensions, domains, ordering, and structural assumptions.
\item Guard small divisors, invalid states, overflow, underflow, and nonfinite values.
\item Make mutation, workspace, indexing, and termination behavior explicit.
\item Test a simple exact case, a difficult valid case, a boundary case, and an expected failure.
\end{itemize}}
\end{cornell}
\begin{summarybox}{Section summary}
Real symmetry guarantees real eigenvalues and an orthonormal eigenbasis, enabling stable reductions and strong residual checks.  Select this technique only when its assumptions match the data, and accept a result only after an independent residual, error estimate, invariant, or refinement check.
\end{summarybox}

\section{11.3 Reduction of a Symmetric Matrix to Tridiagonal Form: Givens and Householder Reductions}
\begin{cornell}
\noterow{What is the central idea?}{Givens or Householder transformations reduce a symmetric matrix to tridiagonal form without changing its eigenvalues.}
\noterow{How should it be used?}{For Reduction of a Symmetric Matrix to Tridiagonal Form: Givens and Householder Reductions, begin from explicit inputs, outputs, preconditions, and representation choices.  Exploit symmetry or sparsity; balance or scale when appropriate; reduce to structured form; iterate and deflate; compute eigenvectors only as needed; verify normalized residuals and orthogonality.}
\noterow{Quantitative lens}{Identify the governing equation, recurrence, probability law, or geometric predicate for Reduction of a Symmetric Matrix to Tridiagonal Form: Givens and Householder Reductions.  Define every variable and scale; then derive a residual, error estimate, or invariant that can be checked independently.}
\noterow{Cost and reliability}{Analyze arithmetic work, storage, data movement, and convergence separately.  Relevant failure pressures include nonnormal sensitivity, clustered eigenvalues, loss of orthogonality, unsafe deflation, incorrect complex-pair handling, and interpreting small residuals without eigenvector conditioning.  Use scaled tolerances and report both success criteria and diagnostic evidence.}
\noterow{Implementation checklist}{\begin{itemize}
\item Validate dimensions, domains, ordering, and structural assumptions.
\item Guard small divisors, invalid states, overflow, underflow, and nonfinite values.
\item Make mutation, workspace, indexing, and termination behavior explicit.
\item Test a simple exact case, a difficult valid case, a boundary case, and an expected failure.
\end{itemize}}
\end{cornell}
\begin{summarybox}{Section summary}
Givens or Householder transformations reduce a symmetric matrix to tridiagonal form without changing its eigenvalues.  Select this technique only when its assumptions match the data, and accept a result only after an independent residual, error estimate, invariant, or refinement check.
\end{summarybox}

\section{11.4 Eigenvalues and Eigenvectors of a Tridiagonal Matrix}
\begin{cornell}
\noterow{What is the central idea?}{Specialized iterations exploit the tridiagonal representation to locate eigenvalues and recover eigenvectors efficiently.}
\noterow{How should it be used?}{For Eigenvalues and Eigenvectors of a Tridiagonal Matrix, begin from explicit inputs, outputs, preconditions, and representation choices.  Exploit symmetry or sparsity; balance or scale when appropriate; reduce to structured form; iterate and deflate; compute eigenvectors only as needed; verify normalized residuals and orthogonality.}
\noterow{Quantitative lens}{Identify the governing equation, recurrence, probability law, or geometric predicate for Eigenvalues and Eigenvectors of a Tridiagonal Matrix.  Define every variable and scale; then derive a residual, error estimate, or invariant that can be checked independently.}
\noterow{Cost and reliability}{Analyze arithmetic work, storage, data movement, and convergence separately.  Relevant failure pressures include nonnormal sensitivity, clustered eigenvalues, loss of orthogonality, unsafe deflation, incorrect complex-pair handling, and interpreting small residuals without eigenvector conditioning.  Use scaled tolerances and report both success criteria and diagnostic evidence.}
\noterow{Implementation checklist}{\begin{itemize}
\item Validate dimensions, domains, ordering, and structural assumptions.
\item Guard small divisors, invalid states, overflow, underflow, and nonfinite values.
\item Make mutation, workspace, indexing, and termination behavior explicit.
\item Test a simple exact case, a difficult valid case, a boundary case, and an expected failure.
\end{itemize}}
\end{cornell}
\begin{summarybox}{Section summary}
Specialized iterations exploit the tridiagonal representation to locate eigenvalues and recover eigenvectors efficiently.  Select this technique only when its assumptions match the data, and accept a result only after an independent residual, error estimate, invariant, or refinement check.
\end{summarybox}

\section{11.5 Hermitian Matrices}
\begin{cornell}
\noterow{What is the central idea?}{Hermitian matrices extend real-symmetric structure to complex data and retain real eigenvalues with unitary eigenvectors.}
\noterow{How should it be used?}{For Hermitian Matrices, begin from explicit inputs, outputs, preconditions, and representation choices.  Exploit symmetry or sparsity; balance or scale when appropriate; reduce to structured form; iterate and deflate; compute eigenvectors only as needed; verify normalized residuals and orthogonality.}
\noterow{Quantitative lens}{Identify the governing equation, recurrence, probability law, or geometric predicate for Hermitian Matrices.  Define every variable and scale; then derive a residual, error estimate, or invariant that can be checked independently.}
\noterow{Cost and reliability}{Analyze arithmetic work, storage, data movement, and convergence separately.  Relevant failure pressures include nonnormal sensitivity, clustered eigenvalues, loss of orthogonality, unsafe deflation, incorrect complex-pair handling, and interpreting small residuals without eigenvector conditioning.  Use scaled tolerances and report both success criteria and diagnostic evidence.}
\noterow{Implementation checklist}{\begin{itemize}
\item Validate dimensions, domains, ordering, and structural assumptions.
\item Guard small divisors, invalid states, overflow, underflow, and nonfinite values.
\item Make mutation, workspace, indexing, and termination behavior explicit.
\item Test a simple exact case, a difficult valid case, a boundary case, and an expected failure.
\end{itemize}}
\end{cornell}
\begin{summarybox}{Section summary}
Hermitian matrices extend real-symmetric structure to complex data and retain real eigenvalues with unitary eigenvectors.  Select this technique only when its assumptions match the data, and accept a result only after an independent residual, error estimate, invariant, or refinement check.
\end{summarybox}

\section{11.6 Real Nonsymmetric Matrices}
\begin{cornell}
\noterow{What is the central idea?}{Nonsymmetric problems may have complex eigenpairs and nonorthogonal eigenvectors, making balancing, Hessenberg reduction, and conditioning important.}
\noterow{How should it be used?}{For Real Nonsymmetric Matrices, begin from explicit inputs, outputs, preconditions, and representation choices.  Exploit symmetry or sparsity; balance or scale when appropriate; reduce to structured form; iterate and deflate; compute eigenvectors only as needed; verify normalized residuals and orthogonality.}
\noterow{Quantitative lens}{Identify the governing equation, recurrence, probability law, or geometric predicate for Real Nonsymmetric Matrices.  Define every variable and scale; then derive a residual, error estimate, or invariant that can be checked independently.}
\noterow{Cost and reliability}{Analyze arithmetic work, storage, data movement, and convergence separately.  Relevant failure pressures include nonnormal sensitivity, clustered eigenvalues, loss of orthogonality, unsafe deflation, incorrect complex-pair handling, and interpreting small residuals without eigenvector conditioning.  Use scaled tolerances and report both success criteria and diagnostic evidence.}
\noterow{Implementation checklist}{\begin{itemize}
\item Validate dimensions, domains, ordering, and structural assumptions.
\item Guard small divisors, invalid states, overflow, underflow, and nonfinite values.
\item Make mutation, workspace, indexing, and termination behavior explicit.
\item Test a simple exact case, a difficult valid case, a boundary case, and an expected failure.
\end{itemize}}
\end{cornell}
\begin{summarybox}{Section summary}
Nonsymmetric problems may have complex eigenpairs and nonorthogonal eigenvectors, making balancing, Hessenberg reduction, and conditioning important.  Select this technique only when its assumptions match the data, and accept a result only after an independent residual, error estimate, invariant, or refinement check.
\end{summarybox}

\section{11.7 The QR Algorithm for Real Hessenberg Matrices}
\begin{cornell}
\noterow{What is the central idea?}{Shifted QR iteration preserves Hessenberg form and drives subdiagonal entries toward deflation, revealing real and complex-conjugate eigenvalues.}
\noterow{How should it be used?}{For The QR Algorithm for Real Hessenberg Matrices, begin from explicit inputs, outputs, preconditions, and representation choices.  Exploit symmetry or sparsity; balance or scale when appropriate; reduce to structured form; iterate and deflate; compute eigenvectors only as needed; verify normalized residuals and orthogonality.}
\noterow{Quantitative lens}{Identify the governing equation, recurrence, probability law, or geometric predicate for The QR Algorithm for Real Hessenberg Matrices.  Define every variable and scale; then derive a residual, error estimate, or invariant that can be checked independently.}
\noterow{Cost and reliability}{Analyze arithmetic work, storage, data movement, and convergence separately.  Relevant failure pressures include nonnormal sensitivity, clustered eigenvalues, loss of orthogonality, unsafe deflation, incorrect complex-pair handling, and interpreting small residuals without eigenvector conditioning.  Use scaled tolerances and report both success criteria and diagnostic evidence.}
\noterow{Implementation checklist}{\begin{itemize}
\item Validate dimensions, domains, ordering, and structural assumptions.
\item Guard small divisors, invalid states, overflow, underflow, and nonfinite values.
\item Make mutation, workspace, indexing, and termination behavior explicit.
\item Test a simple exact case, a difficult valid case, a boundary case, and an expected failure.
\end{itemize}}
\end{cornell}
\begin{summarybox}{Section summary}
Shifted QR iteration preserves Hessenberg form and drives subdiagonal entries toward deflation, revealing real and complex-conjugate eigenvalues.  Select this technique only when its assumptions match the data, and accept a result only after an independent residual, error estimate, invariant, or refinement check.
\end{summarybox}

\section{11.8 Improving Eigenvalues and/or Finding Eigenvectors by Inverse Iteration}
\begin{cornell}
\noterow{What is the central idea?}{Inverse iteration repeatedly solves a shifted linear system, amplifying the eigenvector associated with the eigenvalue nearest the shift.}
\noterow{How should it be used?}{For Improving Eigenvalues and/or Finding Eigenvectors by Inverse Iteration, begin from explicit inputs, outputs, preconditions, and representation choices.  Exploit symmetry or sparsity; balance or scale when appropriate; reduce to structured form; iterate and deflate; compute eigenvectors only as needed; verify normalized residuals and orthogonality.}
\noterow{Quantitative anchor}{\[
(A-\mu I)y_{k+1}=x_k,\qquad x_{k+1}=y_{k+1}/\|y_{k+1}\|
\]
State the dimensions, scaling, and validity assumptions before using this relation.  Check the resulting residual or invariant rather than trusting an iteration count alone.}
\noterow{Cost and reliability}{Analyze arithmetic work, storage, data movement, and convergence separately.  Relevant failure pressures include nonnormal sensitivity, clustered eigenvalues, loss of orthogonality, unsafe deflation, incorrect complex-pair handling, and interpreting small residuals without eigenvector conditioning.  Use scaled tolerances and report both success criteria and diagnostic evidence.}
\noterow{Implementation checklist}{\begin{itemize}
\item Validate dimensions, domains, ordering, and structural assumptions.
\item Guard small divisors, invalid states, overflow, underflow, and nonfinite values.
\item Make mutation, workspace, indexing, and termination behavior explicit.
\item Test a simple exact case, a difficult valid case, a boundary case, and an expected failure.
\end{itemize}}
\end{cornell}
\begin{summarybox}{Section summary}
Inverse iteration repeatedly solves a shifted linear system, amplifying the eigenvector associated with the eigenvalue nearest the shift.  Select this technique only when its assumptions match the data, and accept a result only after an independent residual, error estimate, invariant, or refinement check.
\end{summarybox}

\section*{Method comparison}
\begin{center}
\small
\begin{tabularx}{\linewidth}{>{\bfseries\raggedright\arraybackslash}p{0.22\linewidth}XX}
\toprule
Method or lens & Best use & Main caution \\
\midrule
Jacobi & Small dense symmetric matrices & Simple and accurate; many rotations \\
Tridiagonal plus QR & Dense symmetric eigenproblem & Standard efficient route \\
Inverse iteration & Eigenvector near a known shift & Requires shifted solves \\
\bottomrule
\end{tabularx}
\end{center}

\section*{Worked example}
\begin{exambox}{Independent worked example}
For $A=\begin{bmatrix}2&1\\1&2\end{bmatrix}$, eigenpairs are $(3,(1,1)^T)$ and $(1,(1,-1)^T)$.  Normalize each vector by $1/\sqrt2$ and verify $Av-\lambda v=0$.  Nearly repeated eigenvalues make individual eigenvectors more sensitive even when their invariant subspace is accurate.
\end{exambox}

\section*{Chapter synthesis}
\begin{summarybox}{Chapter synthesis}
The chapter's methods should be viewed as a decision system rather than a list of routines.  Begin with the mathematical problem and its conditioning, exploit structure, select an algorithm with compatible stability and cost, and verify the computed answer with evidence that is independent of the internal iteration.  The recurring implementation discipline is: exploit symmetry or sparsity; balance or scale when appropriate; reduce to structured form; iterate and deflate; compute eigenvectors only as needed; verify normalized residuals and orthogonality.
\end{summarybox}

\subsection*{Common mistakes and numerical hazards}
\begin{warnbox}{Common mistakes and numerical hazards}
Nonnormal sensitivity, clustered eigenvalues, loss of orthogonality, unsafe deflation, incorrect complex-pair handling, and interpreting small residuals without eigenvector conditioning.  A second recurring mistake is reporting only a computed value: always retain tolerances, iteration counts, residuals, refinement behavior, and relevant structural checks.
\end{warnbox}

\subsection*{Retrieval practice}
\textbf{Questions}
\begin{enumerate}
\item What problem does Jacobi Transformations of a Symmetric Matrix address, and which assumption is easiest to violate?
\item What problem does Real Symmetric Matrices address, and which assumption is easiest to violate?
\item What problem does Reduction of a Symmetric Matrix to Tridiagonal Form: Givens and Householder Reductions address, and which assumption is easiest to violate?
\item What problem does The QR Algorithm for Real Hessenberg Matrices address, and which assumption is easiest to violate?
\item What problem does Improving Eigenvalues and/or Finding Eigenvectors by Inverse Iteration address, and which assumption is easiest to violate?
\item How do conditioning, stability, and the final residual play different roles in this chapter?
\end{enumerate}

\textbf{Brief answers}
\begin{enumerate}
\item Jacobi rotations annihilate selected off-diagonal entries while preserving symmetry and eigenvalues, accumulating orthogonal eigenvectors. Recheck the section's domain, structure, scaling, and failure diagnostics.
\item Real symmetry guarantees real eigenvalues and an orthonormal eigenbasis, enabling stable reductions and strong residual checks. Recheck the section's domain, structure, scaling, and failure diagnostics.
\item Givens or Householder transformations reduce a symmetric matrix to tridiagonal form without changing its eigenvalues. Recheck the section's domain, structure, scaling, and failure diagnostics.
\item Shifted QR iteration preserves Hessenberg form and drives subdiagonal entries toward deflation, revealing real and complex-conjugate eigenvalues. Recheck the section's domain, structure, scaling, and failure diagnostics.
\item Inverse iteration repeatedly solves a shifted linear system, amplifying the eigenvector associated with the eigenvalue nearest the shift. Recheck the section's domain, structure, scaling, and failure diagnostics.
\item Conditioning describes sensitivity of the mathematical problem, stability describes error propagation by the algorithm, and the residual measures satisfaction of the computed equations; none is a universal substitute for the others.
\end{enumerate}

\subsection*{Mastery checklist}
\begin{itemize}
\item[$\square$] I can explain and select Jacobi Transformations of a Symmetric Matrix without relying on a memorized code listing.
\item[$\square$] I can explain and select Real Symmetric Matrices without relying on a memorized code listing.
\item[$\square$] I can explain and select Reduction of a Symmetric Matrix to Tridiagonal Form: Givens and Householder Reductions without relying on a memorized code listing.
\item[$\square$] I can explain and select Eigenvalues and Eigenvectors of a Tridiagonal Matrix without relying on a memorized code listing.
\item[$\square$] I can explain and select Hermitian Matrices without relying on a memorized code listing.
\item[$\square$] I can state a scaled stopping rule and an independent verification check.
\item[$\square$] I can identify at least one conditioning risk and one algorithmic stability risk.
\item[$\square$] I can estimate time, storage, and data-movement costs at the appropriate level.
\end{itemize}

\end{document}

